% =====================================================================
% PROJETO DE PESQUISA — PPGEco/UFES
% Proibição da pulverização aérea de agrotóxicos no Ceará (Ensaio 1)
%
% Compilar:  make projeto      (na raiz do repositório)
%       ou:  latexmk -pdf projeto.tex   (dentro de paper/projeto/)
% Motor: pdflatex + bibtex (abntex2cite). Não usar xelatex/biber.
%
% Desenho declarado: docs/ars/20-projeto-de-pesquisa-desenho.md
% Regra do texto: NENHUM número de estimativa. Só números de desenho
% (contagens de municípios, janelas) e magnitudes da literatura.
% Não usar emoji no corpo do texto (quebra o pdflatex).
% =====================================================================
\documentclass[12pt,oneside,a4paper,english,brazil]{abntex2}

\usepackage[T1]{fontenc}
\usepackage[utf8]{inputenc}
\usepackage{lmodern}
\usepackage{indentfirst}
\usepackage{microtype}
\usepackage{amsmath,amssymb}
\usepackage{booktabs}
\usepackage{graphicx}
\usepackage{array}
\usepackage{tikz}
\usetikzlibrary{positioning,arrows.meta,fit}
\usepackage[alf,abnt-emphasize=bf,abnt-etal-list=0]{abntex2cite}

% --- Dados do trabalho (edite aqui) ----------------------------------
\titulo{Proibição da pulverização aérea de agrotóxicos no Ceará:
        efeitos sobre o peso ao nascer}
\autor{Felipe Carvalho Souza Santos}
\orientador{Prof. Dr. Renato Seixas}
\instituicao{%
  Universidade Federal do Espírito Santo
  \par Centro de Ciências Jurídicas e Econômicas
  \par Programa de Pós-Graduação em Economia}
\local{Vitória}
\data{2026}
\tipotrabalho{Projeto de pesquisa}
\preambulo{Projeto de pesquisa do primeiro ensaio da dissertação, apresentado
ao Programa de Pós-Graduação em Economia da Universidade Federal do Espírito
Santo, como requisito parcial do curso de Mestrado em Economia.}

% --- Ambiente "quadro" (ABNT distingue quadro de tabela) ------------------
\newcommand{\quadroname}{Quadro}
\newfloat{quadro}{loq}{\quadroname}
\counterwithout{quadro}{chapter}
\counterwithout{equation}{chapter}

% --- Aparência ---------------------------------------------------------
\hypersetup{colorlinks=true,linkcolor=black,citecolor=black,urlcolor=black,
  pdftitle={Proibição da pulverização aérea de agrotóxicos no Ceará: efeitos sobre o peso ao nascer},pdfauthor={Felipe Carvalho Souza Santos}}
\makeatletter\g@addto@macro{\UrlBreaks}{\do\-}\makeatother % URLs longas quebram em hífen
\setlength{\parindent}{1.25cm}
\setlength{\parskip}{0.1cm}
% Projeto curto: capítulos em sequência, sem quebra de página entre eles.
\renewcommand{\clearforchapter}{}
\setlength{\beforechapskip}{12pt}
\setlength{\afterchapskip}{12pt}

\begin{document}
\selectlanguage{brazil}
\frenchspacing

% Capa: instituição, autor, título, local e ano (NBR 14724)
\begin{capa}
  \center
  {\large\textsc{Universidade Federal do Espírito Santo}\par
   \textsc{Centro de Ciências Jurídicas e Econômicas}\par
   \textsc{Programa de Pós-Graduação em Economia}\par}
  \vspace*{3cm}
  {\ABNTEXchapterfont\large\imprimirautor}
  \vfill
  {\ABNTEXchapterfont\bfseries\LARGE\imprimirtitulo}
  \vfill
  {\large Projeto de pesquisa}
  \vfill
  {\large\imprimirlocal\par\imprimirdata}
  \vspace*{1cm}
\end{capa}
\imprimirfolhaderosto

% Resumo (pré-textual)
\begin{resumo}
Em 2019, o Ceará proibiu a pulverização aérea de agrotóxicos em todo o seu
território. Este projeto, que trata do primeiro ensaio da dissertação, propõe
avaliar o efeito dessa proibição sobre o peso ao nascer das coortes cujo
início de gestação cairia na estação de pulverização, em relação às expostas
à estação apenas no fim da gestação. O desenho principal é uma diferença
tripla de calendário gestacional: compara essas coortes, antes e depois da
lei, em municípios com aplicação aérea documentada entre 2007 e 2011 e em
municípios de bananicultura familiar. O período de referência é fixado antes
da seca de 2012--2017, que pode ter suspendido os voos e que afeta o peso ao
nascer por vias próprias; ler o resultado como efeito da lei supõe que, sem
ela, os voos teriam voltado depois da seca. O projeto trata como problema
central a distância entre a cultura plantada e o método de aplicação, propõe
medi-la com registros administrativos e fixa de antemão o que entrega em cada
cenário, inclusive o de o teste não se realizar por falta de registros ou de
poder. Duas rodadas exploratórias já cruzaram a exposição com o desfecho, e o
texto as declara.

\textbf{Palavras-chave}: agrotóxicos; pulverização aérea; peso ao nascer;
avaliação de impacto; Ceará.
\end{resumo}

\pdfbookmark[0]{\contentsname}{toc}
\tableofcontents*
\clearpage

\textual
% Capítulo 1 — Introdução. Reescrita em 2026-11-16 pela fórmula de Head (via
% Bellemare, Doing Economics, cap. 2; docs/framework-dissertacao.md §3.1):
% gancho no mundo real, pergunta como primeira frase, antecedentes como
% história, contribuição, organização. Promete só o que o corpo entrega.

\chapter{Introdução}
\label{cap:introducao}

Em abril de 2010, José Maria Filho, o Zé Maria do Tomé, foi assassinado na
Chapada do Apodi, no interior do Ceará. Ele denunciava havia anos a
pulverização aérea de agrotóxicos sobre os perímetros irrigados da região, onde
aviões sobrevoavam plantações de banana e melão vizinhas a casas, escolas e
açudes de abastecimento. A denúncia não tratava do veneno aplicado na lavoura,
mas do veneno que não fica na lavoura. Nove anos depois, o Ceará proibiu a
prática em todo o território e deu à lei o nome do homem assassinado (Lei
Estadual nº~16.820/2019).

A pulverização aérea é um caso de externalidade em que o mecanismo físico é
visível. Parte da calda liberada pela aeronave não atinge o alvo: é carregada
pelo vento e se deposita fora da área tratada, fenômeno que a literatura
técnica chama de deriva. Quem contrata a aplicação paga pelo serviço e pelo
produto; quem mora a favor do vento arca com o restante. A lei cearense
responde a esse conflito por proibição, e não se localizou, na literatura
consultada, avaliação causal dos efeitos dessa resposta.

\textbf{Qual é o efeito causal da proibição da pulverização aérea de
agrotóxicos sobre o peso ao nascer nos municípios cearenses onde havia
aplicação aérea?} É a pergunta deste projeto, e ela é de avaliação de política
pública. O objeto é a retirada de um \emph{método de aplicação}, e não de uma
substância, e essa distinção organiza o desenho inteiro: o grupo tratado
relevante é o dos municípios onde havia avião, e não o dos municípios onde
havia lavoura. A resposta buscada é um contraste, nos municípios com aplicação
aérea documentada entre 2007 e 2011: o efeito sobre as coortes cujo início de
gestação cairia na estação de pulverização, em relação às expostas a ela
apenas no fim da gestação. A estação de pulverização são os meses em que os voos se
concentram, provisoriamente a quadra chuvosa, de fevereiro a maio; o período
de referência, as coortes nascidas de 2007 a 2011.

\section{Antecedentes na literatura}

A evidência econômica próxima se divide entre o que uma substância faz à saúde
e o que um método de aplicação faz à saúde. Do lado da substância,
\citeonline{dias2023down} exploram a expansão da soja transgênica no Brasil,
combinada à direção do fluxo d'água dentro das bacias, e documentam piora dos
desfechos ao nascer nas populações a jusante, e
\citeonline{reynier2025glyphosate} associam a difusão do glifosato nos Estados
Unidos a menor peso ao nascer e gestação mais curta.

Do lado do método, os antecedentes diretos são dois.
\citeonline{camacho2017health} avaliam as consequências sanitárias da aspersão
aérea de herbicida sobre cultivos ilícitos na Colômbia.
\citeonline{calzada2023bananas} estudam a fumigação aérea de bananais no
Equador e encontram déficit de peso ao nascer entre filhos de mães que moravam
junto às plantações e foram expostas, na gestação, aos períodos de fumigação
intensa. Este último é o análogo mais próximo do caso cearense, pela cultura e
pelo método, com diferenças que importam: a doença combatida não é a mesma e,
pelos relatos disponíveis, a frequência das aplicações é menor
(Capítulo~\ref{cap:contexto}).

Entre os trabalhos localizados, nenhum avalia uma \textbf{proibição de
método}. Há simulações do banimento de moléculas específicas
\cite{bocker2020glyphosate}, mas simulação \emph{ex ante} não é avaliação
\emph{ex post}. Nos estudos de método, a pulverização aérea é introduzida ou
está em curso; no Ceará ela é retirada, e os sinais esperados se invertem.

\section{Justificativa}

A questão tem relevância prática imediata. A lei cearense foi contestada no
Supremo Tribunal Federal, que julgou constitucional a vedação em 2023
\cite{stf2023adi6137}, o que assentou a competência dos estados para legislar
sobre a matéria e torna a medida replicável; em 2024 o próprio Ceará
flexibilizou a regra para admitir drones, e a nova redação voltou a ser
questionada na Corte \cite{stf2025adi7794}. O debate tem se apoiado em
argumentos de plausibilidade, porque não se localizou estimativa causal do que
a proibição de fato produziu. Uma avaliação que diga se houve ganho em peso ao
nascer, ou que mostre com rigor por que ele não é mensurável com os registros
existentes, informa os estados que consideram adotar a medida e a decisão de
mantê-la, relaxá-la ou substituí-la. O alcance é limitado: o ensaio informa um
desfecho, por uma via, em nível municipal, e não mede os demais danos
atribuídos à pulverização.

\section{Contribuição e organização do projeto}

A contribuição pretendida tem três camadas. A primeira é de \textbf{objeto}: a
avaliação \emph{ex post} de uma proibição de método de aplicação. A segunda é
de \textbf{medida}: o projeto trata como problema central a distância entre a
cultura plantada e o método de aplicação efetivamente usado, e propõe medi-la
com registros administrativos; ela depende de os registros serem obtidos. A
terceira é de \textbf{identificação}: como a lei vale para todo o estado desde
a mesma data, o projeto usa um desenho de calendário gestacional
\cite{calzada2023bananas,brainerd2014seasonal}, com os municípios de banana
familiar como grupo de comparação e um período de referência fixado de
antemão, e o confronta com um desenho de intensidade de exposição
\cite{callaway2024continuous}. O teste só se realiza se houver voos
documentados e poder; em todos os cenários previstos, inclusive o de ele não
se realizar, o projeto diz o que entrega (Capítulo~\ref{cap:resultados}).

Este documento é o projeto do primeiro ensaio da dissertação; o segundo, sobre
a escolha do instrumento, é apenas mencionado e terá projeto próprio.
Descreve-se um plano de trabalho, não um resultado, embora duas rodadas
exploratórias já tenham cruzado a exposição com o desfecho do pós-lei
(Seção~\ref{sec:inferencia}). O Capítulo~\ref{cap:contexto} apresenta a lei e
a região; o Capítulo~\ref{cap:referencial}, a literatura; o
Capítulo~\ref{cap:objetivos}, os objetivos e as hipóteses; o
Capítulo~\ref{cap:metodologia}, a estratégia empírica; o
Capítulo~\ref{cap:dados}, os dados; e os
Capítulos~\ref{cap:resultados} e~\ref{cap:cronograma}, os resultados esperados
e o cronograma.

% Seção 2 — Contexto institucional. Condensada de paper/secoes/02-background.tex
% e docs/legislacao/README.md. Fatos de 2011: docs/ars/19 §2.3 (a conferir).

\chapter{Contexto institucional}
\label{cap:contexto}

\section{A Lei Zé Maria do Tomé}

A Lei Estadual nº~16.820, sancionada em 8 de janeiro de 2019 e vigente desde o
dia seguinte \cite{ceara2019lei16820}, acrescentou o art.~28-B à Lei Estadual
nº~12.228/1993 \cite{ceara1993lei12228}, a norma cearense sobre agrotóxicos. O caput veda ``a pulverização aérea de agrotóxicos
na agricultura no Estado do Ceará''; o §~2º estende a vedação à dispersão por
aeronave para controle vetorial. Dois traços da redação importam ao desenho. A
proibição é de \textbf{método}: nenhum princípio ativo é nomeado, e o produto
continua disponível para aplicação terrestre. E a proibição é \textbf{estadual
e simultânea}: alcança os 184 municípios na mesma data.

O Quadro~\ref{qua:cronologia} resume a cronologia. Três marcos antecedem o
tratamento e são tratados separadamente: a \emph{notícia} (apresentação do projeto de lei que deu origem à norma, em
2015, precedida de proposição semelhante em 2011), a \emph{certeza} (aprovação unânime, em dezembro de
2018) e a \emph{obrigação} (vigência, em janeiro de 2019). Em dezembro de 2024,
a Lei nº~19.135 \cite{ceara2024lei19135} deu nova redação ao artigo e abriu exceção para aeronaves
remotamente pilotadas. O estado não revogou a proibição, mas trocou o
instrumento: de cota zero para regulação de desempenho. Por isso a janela de
análise termina em 19 de dezembro de 2024.

\begin{quadro}[htb]
\caption{\label{qua:cronologia}Cronologia legal do tratamento}
\centering
\small
\begin{tabular}{@{}lp{11.2cm}@{}}
\toprule
Data & Evento \\
\midrule
09/12/1993 & Lei Estadual nº 12.228: norma-mãe sobre agrotóxicos no Ceará. \\
20/11/2009 & Limoeiro do Norte proíbe a pulverização aérea por lei municipal (Lei nº 1.478). \\
2010 & Assassinato de Zé Maria do Tomé (abril). Em maio, a lei municipal é revogada (Lei nº 1.511). \\
2011 & Notícia de proibição da prática na Chapada do Apodi (março); ação civil pública do Ministério Público e projeto de lei estadual com o mesmo objeto (julho). Alcance, desfecho e tramitação a conferir. \\
24/02/2015 & Apresentação do Projeto de Lei nº 18/2015 (marco de notícia). \\
18/12/2018 & Aprovação por unanimidade na Assembleia Legislativa (marco de certeza). \\
09/01/2019 & Vigência da Lei nº 16.820 (marco de obrigação). \\
05/2023 & O Supremo Tribunal Federal julga a lei constitucional (ADI 6137). \\
19/12/2024 & Lei nº 19.135: exceção para drones. Fim da janela de análise. \\
\bottomrule
\end{tabular}
\fonte{Elaboração própria, a partir dos textos legais e da tramitação legislativa; os eventos de 2009 a 2011 provêm de fontes secundárias, a conferir.}
\end{quadro}

\section{A Chapada do Apodi, a cultura e o químico}

A Chapada do Apodi, no Baixo Jaguaribe, concentra a fruticultura irrigada do
Ceará, com destaque para banana e melão. É a região à qual a lei responde e
onde a exposição da população à deriva foi documentada antes da proibição. Na
bananicultura, a aplicação aérea serve ao controle de doenças fúngicas da
folha, com aplicações repetidas ao longo do ano. No análogo equatoriano, elas
se intensificam na estação chuvosa, quando a pressão da doença aumenta
\cite{calzada2023bananas}; se o mesmo vale para o Ceará é questão que a etapa
de medição deste projeto verifica.

Esse perfil distingue o caso da literatura de referência, construída em grande
parte sobre glifosato em grãos transgênicos. A tabela de análises de água
reproduzida na justificativa do Projeto de Lei nº~18/2015, atribuída ao Dossiê
ABRASCO \cite{ceara2015pl18}, registra 23 pontos de coleta na Chapada em 2009: o fungicida
procimidona e o inseticida carbaril aparecem em todos, carbofurano em 18,
fenitrotiona em 16 e glifosato em apenas 4. Trata-se de fungicida e inseticida
sobre fruta irrigada. Os desenhos de referência transferem na forma, não no
mecanismo, e o análogo empírico adequado é o da fumigação aérea de bananais
\cite{calzada2023bananas}.

\section{O que se sabe sobre onde havia avião}
\label{sec:onde-aviao}

O dado público mais direto sobre o uso de aeronaves é o Censo Agropecuário de
2006, que registra os estabelecimentos que aplicaram agrotóxicos por avião. No Ceará eram 36 estabelecimentos, em 7 dos 184 municípios
\cite{ibge2006censoagro}. A prática era, portanto, geograficamente concentrada,
ao passo que a bananicultura está presente na maior parte do estado: a maioria
dos municípios com a cultura não tinha registro de avião.

Dois fatos posteriores ao Censo impedem tomar esse retrato como definitivo. Em
março de 2011, noticiou-se que as empresas de fruticultura da Chapada estariam
proibidas de pulverizar por aeronave, e em julho do mesmo ano o Ministério
Público ajuizou ação civil pública com esse pedido. O alcance da primeira
medida e o desfecho da segunda não foram ainda conferidos em fonte primária.
Saber se a prática persistiu até 2018 é condição para saber o que a lei de
2019 removeu, e por isso a etapa de medição do
Capítulo~\ref{cap:metodologia} antecede a estimação principal.

% Seção 3 — Referencial teórico e empírico. Reescrita em 2026-11-15: referencial
% de instrumento reduzido a um parágrafo; Calzada et al. corrigido pelo texto
% aceito (doc de 2026-11-10, §4). Magnitudes só da literatura.

\chapter{Referencial teórico e empírico}
\label{cap:referencial}

\section{A deriva como externalidade negativa}

O agrotóxico que deriva sobre o vizinho é um caso clássico de externalidade
negativa: o produtor que contrata a aplicação aérea internaliza o custo do
serviço e do produto, mas não o custo em saúde suportado por quem mora a favor
do vento. A solução negociada de \citeonline{coase1960} supõe custos de
transação baixos; aqui a população atingida é difusa e o dano é
probabilístico e de difícil atribuição, o que justifica a intervenção
regulatória. Qual instrumento usar, se proibição, taxa ou regulação de
desempenho, depende das inclinações relativas do benefício e do custo
marginais \cite{weitzman1974,skevas2012} e tem dimensão distributiva
\cite{lichtenberg1988,helfand1995}. Essa escolha é objeto de um
segundo ensaio da dissertação. O que este projeto fornece a ela é a premissa
empírica: saber se a proibição produziu o benefício em saúde que a motivou.

\section{Agrotóxicos e saúde ao nascer}

A literatura econômica recente sobre exposição gestacional a agrotóxicos
converge em três regularidades úteis ao desenho.

A primeira é a ordem de grandeza e a sua dependência da proximidade.
\citeonline{reynier2025glyphosate} estimam reduções médias de peso ao nascer da
ordem de 23 a 32 gramas associadas à difusão do glifosato em áreas rurais dos
Estados Unidos. \citeonline{calzada2023bananas} encontram déficit de cerca de
80 gramas, mas apenas para mães que moravam a até 150 metros do perímetro de
um bananal e tiveram três meses ou mais de fumigação intensiva durante a
gestação; para as que moravam mais longe não há gradiente. No Brasil,
\citeonline{rigotto2013trends} comparam municípios da fruticultura irrigada do
Ceará com municípios de agricultura familiar e documentam piores indicadores
de saúde nos primeiros, em estudo ecológico, sem pretensão causal.

A segunda é a concentração. Os efeitos aparecem no topo da distribuição de
exposição \cite{larsen2017agricultural} e na cauda inferior da distribuição de
peso esperado \cite{reynier2025glyphosate}; entre os poucos estudos que isolam
fungicidas está o de \citeonline{jones2020invasive}. A consequência para este
projeto é direta: se o efeito existe só para quem mora junto ao bananal, a
média de um município inteiro o dilui pela fração de nascimentos exposta, e um
desenho municipal pode não encontrá-lo mesmo que ele seja grande.

A terceira é o calendário. Quando a aplicação tem estação, o efeito é
identificado pelo cruzamento entre o período sensível da gestação e a estação
de aplicação: \citeonline{calzada2023bananas} encontram efeito nos dois
primeiros trimestres e estimativa sem significância estatística no terceiro, com a intensidade da fumigação
ditada pela estação chuvosa, e \citeonline{brainerd2014seasonal} usam o mês da
concepção e a estação de aplicação de agroquímicos na Índia. A ameaça
conhecida é a composição sazonal das mães: quem concebe em cada época do ano
difere em características que afetam o desfecho
\cite{buckles2013season,currie2013within}.

\section{Avaliação de uma política estadual e simultânea}

Como a lei vale para todo o estado desde a mesma data, os estimadores e as
decomposições de diferenças-em-diferenças para adoção escalonada
\cite{callaway2021multiple,goodmanbacon2021,dechaisemartin2020} não se aplicam:
não há variação de \emph{timing}. Restam a variação de \emph{calendário} entre
coortes de nascimento, para a qual a diferença tripla exige uma única hipótese
de tendências paralelas, a da diferença entre os dois contrastes que a compõem
\cite{olden2022triple}, e a variação de \emph{intensidade} prévia entre
municípios, para a qual \citeonline{callaway2024continuous} mostram que o
efeito no nível da dose recebida sai sob tendências paralelas usuais, ao passo
que a curva dose-resposta exige hipótese mais forte.

Quando poucos grupos são tratados, a inferência assintótica usual falha e são
necessários procedimentos próprios \cite{conley2011,ferman2019inference,
chernozhukov2021conformal,hagemann2025single}.

% Seção 4 — Objetivos e hipóteses. Reescrita em 2026-11-15.

\chapter{Objetivos e hipóteses}
\label{cap:objetivos}

\textbf{Objetivo geral.} Estimar o efeito causal da proibição estadual da
pulverização aérea de agrotóxicos (Lei nº~16.820/2019) sobre a saúde ao nascer
nos municípios cearenses onde a prática existia.

\textbf{Objetivos específicos.}
\begin{alineas}
  \item medir onde e quando havia aplicação aérea antes da proibição, com
    registros administrativos federais e estaduais;
  \item verificar, só com dados anteriores à lei, se a exposição à estação de
    pulverização deixava marca sazonal no peso ao nascer;
  \item estimar o efeito da proibição pelo desenho de calendário gestacional,
    com os municípios de banana familiar como grupo de comparação;
  \item confrontar essa estimativa com a forma reduzida estadual, o controle
    sintético e a comparação de fronteira;
  \item testar a seleção para o nascimento vivo e o canal de substituição da
    aplicação aérea pela terrestre.
\end{alineas}

\textbf{Hipóteses.} A hipótese principal ($H_1$) é que, após a lei, o peso ao
nascer das coortes gestadas na estação de pulverização aumentou, em relação às
demais coortes, nos municípios onde havia aplicação aérea, comparados aos de
banana familiar. A segunda hipótese confirmatória ($H_2$) é que a taxa de
óbito fetal dessas coortes caiu. Como hipóteses de mecanismo, o efeito deve
ser maior nos anos chuvosos ($H_3$) e ausente quando só o terceiro trimestre
coincide com a estação e nos municípios de banana familiar comparados aos sem
banana ($H_4$). Como hipótese de substituição ($H_5$), as notificações de
intoxicação ocupacional por agrotóxico agrícola podem aumentar, se a aplicação
migrou do avião para o trator e o pulverizador costal.

% Seção 5 — Metodologia. Revisão do Estágio 4 (2026-11-15), após a decisão
% "revisão maior" do painel (docs/ars/pipeline-rodada20/stage3/). Decisões do
% pesquisador: reparar o desenho de calendário; diagnóstico descritivo, sem
% poder de decisão; pré-período de referência fixado (2007–2011); peso ao
% nascer como única hipótese confirmatória. NENHUM número de estimativa aqui.
% Revisão da rodada 21 (2026-11-16), após a decisão do Estágio 3: parâmetro
% como contraste, hipótese de retorno dos voos, regra de decisão completada
% no Apêndice F. Identificação e estimador principal não mudaram.

\chapter{Metodologia}
\label{cap:metodologia}

\section{Parâmetro-alvo e cadeia causal}

O parâmetro de interesse é o efeito da lei, em intenção de tratar, sobre a
diferença de peso ao nascer entre as coortes cujo início de gestação cai na
estação de pulverização e as expostas a ela apenas no fim da gestação, nos
municípios cearenses com aplicação aérea documentada no período de referência.
Com $\Delta_{mt} = Y_{mt}(1) - Y_{mt}(0)$ a diferença entre os desfechos com e
sem a proibição, no pós-lei,
\begin{equation}
  \tau = E\left[\,\Delta_{mt} \mid A_m = 1,\ H_t = 1\,\right]
       - E\left[\,\Delta_{mt} \mid A_m = 1,\ R_t = 1\,\right],
\end{equation}
em que $A_m$ indica município com aplicação aérea documentada entre 2007 e
2011, $H_t$, coorte com exposição alta à estação no início da gestação, e
$R_t$, coorte de referência, exposta apenas no terceiro trimestre; as médias
são por nascimento. O contraste só limita por baixo o efeito em nível sob uma
hipótese mantida, não testável: a de que a exposição tardia não beneficia o
feto. O tratamento é a lei, e não o voo: o que se estima inclui o
descumprimento e a substituição pela aplicação terrestre. A referência de
$A_m$ é anterior à seca de 2012--2017 \cite{rebello2018seca}, em que os voos
podem ter sido suspensos (Seção~\ref{sec:onde-aviao}); definir o tratado pelo
que ocorria às vésperas da lei poderia deixar o grupo vazio. Essa escolha
exige uma hipótese de identificação adicional, não verificada: sem a lei, os
voos teriam voltado ao padrão da referência depois da seca. Sob ela, $\tau$ é
o efeito de a lei ter impedido os voos onde eles ocorreriam. Os relatórios de
2008 a 2024 e o relato de voo até 2018 são a evidência que a examina; se os
voos cessaram antes da lei por outra causa, um coeficiente positivo não
distingue o efeito da lei do fim da prática (Quadro~\ref{qua:regra}, linha 5).

A cadeia que liga a lei ao desfecho tem elos que podem falhar separadamente: a
lei precisa ter sido cumprida; precisava haver voos a suspender ou a impedir
de voltar; a ausência de voos precisa reduzir a exposição das gestantes, e não
apenas trocá-la pela da aplicação terrestre; e a menor exposição precisa
alterar o desfecho. Um efeito nulo no desfecho final não distingue falha de implementação de
falha do mecanismo, e por isso o desenho mede os elos intermediários.

\section{Desenho principal: diferença tripla de calendário}
\label{sec:calendario}

O Apêndice~\ref{ap:dag} traz o diagrama causal do desenho e liga cada
elemento dele à hipótese de identificação correspondente. Dentro de cada município, comparam-se as coortes de
exposição alta à estação no início da gestação com as expostas apenas no fim;
essa diferença, o hiato sazonal, é comparada entre períodos e entre os
municípios com aplicação aérea e os de comparação.

Seja $P_t \in [0,1]$ a fração dos dois primeiros trimestres de gestação da
coorte $t$ que cai na estação de pulverização (convenções no
Apêndice~\ref{ap:espec}). As coortes são classificadas em
três grupos, fixados neste projeto: exposição alta, $H_t = 1$, se $P_t$ é de
metade ou mais; intermediária, $I_t = 1$, se $P_t$ é positivo e menor que a
metade; e nula no início da gestação, a referência, se $P_t = 0$. Sejam ainda
$Seca_t$ o indicador das coortes nascidas de 2013 a 2018, cujo início de
gestação cai, em geral, nos anos de seca, $Pos_t$ o das
nascidas a partir de 2020 e $C_{m,q(t)}$ a chuva, no município, da quadra
chuvosa a que o início da gestação foi exposto. Estima-se, com ponderação pelo número de nascimentos,
\begin{equation}
\begin{split}
  Y_{mt} ={}& \alpha_{m,s(t)} + \mu_{m,a(t)} + \gamma_t
   + \sum_{k \in \{I,H\}} \Bigl[\theta_k\,(k_t \times C_{m,q(t)})
   + \beta^k_S\,(A_m \times k_t \times Seca_t) \\
  & + \beta^k\,(A_m \times k_t \times Pos_t)\Bigr] + \varepsilon_{mt},
\end{split}
\label{eq:ddd}
\end{equation}
em que $\alpha_{m,s(t)}$ são efeitos fixos de município por mês do ano,
$\mu_{m,a(t)}$ de município por ano de nascimento e $\gamma_t$ de coorte. Os
primeiros absorvem a sazonalidade de base de cada município; os segundos, os
choques comuns a todas as coortes de um município num ano, entre eles os
programas estaduais de saúde materna contemporâneos à lei; as coortes de
transição, nascidas em 2012 e em 2019, são excluídas. Os termos de
ordem inferior são absorvidos pelos efeitos fixos. O coeficiente de
interesse é $\beta \equiv \beta^H$: quanto a diferença de peso ao nascer entre
as coortes de exposição alta e as de referência mudou, do período de
referência para o pós-lei, nos municípios com aplicação aérea em relação aos
de comparação. Sob a hipótese principal, $\beta > 0$. O coeficiente
$\beta_S \equiv \beta^H_S$ mede a mesma variação no período de seca e é
descritivo.

A exposição entra em grupos, e não como variável contínua, porque o
coeficiente de uma exposição contínua mistura efeitos de níveis diferentes,
por analogia com \citeonline{callaway2024continuous}
(Capítulo~\ref{cap:referencial}). Com uma
estação de quatro meses, toda gestação a termo a atravessa em algum trimestre:
as coortes de referência têm apenas o terceiro na estação, e por isso o
parâmetro é um contraste, que subestima o efeito em nível se a exposição
tardia também causar dano. A janela gestacional é fixa: a duração observada da
gestação é desfecho, e usá-la para definir a exposição condicionaria em
variável afetada pelo tratamento.

\textbf{Período de referência.} São as coortes nascidas de 2007 a 2011, anos
em que há relatos de voos na Chapada (Seção~\ref{sec:onde-aviao}) e que
antecedem a seca. A escolha foi registrada em 15 de novembro de 2026, depois
da rodada exploratória de setembro e antes da de novembro, que a tomou como
dada (Apêndice~\ref{ap:regra}). As coortes de 2005 e 2006 entram apenas
no estudo de evento. Três limites são conhecidos: os registros pedidos começam
em 2008; em 2011 houve projeto de lei estadual e ação civil pública pela
proibição; e Limoeiro do Norte teve proibição municipal de novembro de 2009 a
maio de 2010. Reportam-se a referência restrita a 2008--2010 e a exclusão das
coortes de Limoeiro afetadas (Apêndice~\ref{ap:espec}).

\textbf{Grupo de comparação.} Os municípios de comparação são os produtores de
banana em que a cultura é predominantemente familiar, segundo o Censo
Agropecuário de 2017 \cite{ibge2017censoagro}, com limiar a fixar no plano de
análise. Supõe-se que não tenham aplicação aérea, o que se verifica nos
registros e no Censo de 2006, que a registra também fora dos polos irrigados;
município com qualquer registro de aeronave é excluído. Em nível, a comparação
não seria válida: agricultura familiar e empresarial diferem em renda, acesso
a programas sociais e exposição à aplicação costal. Dentro da diferença
tripla exige-se menos: que, na ausência da lei, o hiato sazonal tivesse
evoluído do mesmo modo nos dois grupos desde o período de referência
\cite{olden2022triple}. A hipótese é examinada por um estudo de evento,
com um coeficiente por ano e 2011 como base, que fora da referência cobre só
2005 e 2006, com pouco poder, e por um tratamento falso, que compara os
municípios de banana familiar aos sem banana e examina o grupo de comparação,
não o contrafactual dos tratados. Os dois são relatados junto a $\beta$, e não
como portão: um resultado desfavorável não suspende o teste, e $\beta$ é
sempre lido com intervalo de sensibilidade a desvios das tendências paralelas
do tamanho observado antes da lei
\cite{rambachan2023credible}. Como a identificação pode depender de características
dos municípios, reporta-se, se o número de tratados permitir, um estimador
duplamente robusto com covariáveis anteriores à lei \cite{ortiz2025triple}.

\textbf{Chuva.} A chuva entra de dois modos. Como a precipitação durante a
gestação afeta o peso ao nascer no semiárido por vias que não passam pelo
avião \cite{rocha2015water}, a chuva da quadra chuvosa, interagida com os indicadores de exposição, é controle nos dois grupos, sob a hipótese adicional, declarada, de que esse efeito
direto seja comum a municípios de lavoura irrigada e de sequeiro. É ela que
separa a lei do retorno das chuvas; por isso a conclusão confirmatória exige
que a especificação principal e a que deixa o efeito diferir por grupo
concordem em sinal e em ordem de grandeza (Apêndice~\ref{ap:regra}), e, se
divergirem, $\beta$ é reportado como resultado não separado do retorno das
chuvas. E, como a doença que motiva o
voo depende de umidade, examina-se, em caráter exploratório, se $\beta$ é maior
nos anos chuvosos. As especificações estão no Apêndice~\ref{ap:espec}.

\textbf{Diagnóstico descritivo do pré-período.}\label{sec:diagnostico} Antes
da análise confirmatória, descreve-se o hiato sazonal nos dois grupos no
período de referência e no período de seca. O diagnóstico responde se a
exposição à estação aparece nos dados municipais e se a seca a alterou. Ele é
descritivo: não escolhe o período de referência e não autoriza nem impede a
análise confirmatória (Seção~\ref{sec:regra}; a razão está no
Apêndice~\ref{ap:ameacas}). A leitura tem um limite: os dois períodos diferem
pela própria seca, que altera a sazonalidade do peso ao nascer por outras
vias, de modo que o fechamento do hiato nas coortes de 2013 a 2018 é
compatível com a suspensão dos voos, mas não a demonstra.

\section{Medição do método}
\label{sec:medicao}

A classificação $A_m$ e o calendário de pulverização são medidos, e não
supostos. A fonte principal são os relatórios mensais que as empresas de
aviação agrícola entregam ao Ministério da Agricultura \cite{mapa2008in2}, com
município, cultura, área e produto; pedem-se os de 2008 a 2024, para o Ceará e
o Rio Grande do Norte, incluídos os operadores sediados em outros estados, e a
relação de operadores registrados, com a qual se afere a completude da série.
Os anos posteriores à lei medem o cumprimento e, desde 2021, o uso de drones
\cite{mapa2021portaria298}. Não se sabe se os relatórios anteriores ao envio
eletrônico foram preservados, e por isso a entrega da medida é condicional; o
silêncio de uma série incompleta não é lido como ausência de voo (regra de
classificação e critério de completude no Apêndice~\ref{ap:regra}).
Complementam a medida o cadastro estadual de prestadores, os autos de
infração, o receituário agronômico e os autos da ação civil pública de 2011
(Apêndice~\ref{ap:fontes}). O Censo Agropecuário de 2006
\cite{ibge2006censoagro} fornece a classificação provisória; por não
distinguir cultura e ser anterior à seca, pode apontar para municípios
diferentes daqueles em que se voava em 2018.

\section{Interferência e versões do tratamento}

O desenho supõe que o tratamento de um município não altera o desfecho de
outro e que o tratamento tem uma só versão. Quatro canais de interferência
são conhecidos: a deriva sobre municípios vizinhos e o aquífero compartilhado,
que atenuariam a estimativa; a migração do serviço aéreo para áreas sem
proibição, que a inflaria; e a mobilidade das mães. Como o análogo equatoriano não encontra efeito além de 150 metros do
bananal \cite{calzada2023bananas}, a interferência por deriva deve ser
pequena, mas os vizinhos imediatos dos tratados são separados dos demais,
formando três grupos. Quanto às versões, estar sob a lei significa coisas
diferentes: ter parado de voar; ter trocado o avião pela aplicação terrestre
ou pelo drone; ou não ter cumprido. O parâmetro, de intenção de tratar, é a
média sobre essas três versões. Não voar desde antes da lei não é uma quarta
versão: onde os voos cessaram antes de 2019, pela seca, pela regulação
municipal ou por decisão judicial, a lei não alterou a exposição, e $\beta$
atribuiria a ela o que tem outra causa.

\section{Análises secundárias e triangulação}
\label{sec:secundarias}

Três análises complementam o desenho principal. A primeira é a forma reduzida
estadual: um diferenças-em-diferenças com tratamento contínuo
\cite{callaway2024continuous}, em que a dose é a área plantada de banana antes
da lei \cite{ibgepam}, o grupo de dose nula é construído também pela aptidão
agroclimática e o alvo é o efeito no nível da dose recebida. Como a área mede
a cultura e não o método, classificar pela cultura produz falsos positivos, que
atenuam a comparação \cite{denteh2022misclassification}, e o Censo de 2006
sugere que há também falsos negativos; por isso a forma reduzida é lida como
descritiva. A segunda é o controle sintético e o
diferenças-em-diferenças sintético \cite{arkhangelsky2021sdid} para o agregado
dos municípios com avião, com doadores fora do alcance da interferência. A
terceira é a comparação com os municípios potiguares da Chapada do Apodi, que
não é experimento natural limpo, porque tudo o que é estadual muda junto na
divisa e as culturas diferem. A comparação entre as análises é descritiva, sem
rótulo de concordância: cada uma tem estimando próprio, e reporta-se o valor
que se esperaria de cada uma sob o efeito principal.

\section{Regra de decisão}
\label{sec:regra}

O Quadro~\ref{qua:regra} fixa o que cada resultado da medição e do cálculo de
poder aciona; o Apêndice~\ref{ap:regra} traz definições, critérios numéricos
e o enquadramento dos casos intermediários, entre eles os registros parciais
e os tardios. Os critérios foram escritos depois das rodadas exploratórias,
com os resultados delas à vista; nenhum é calculado com o desfecho do pós-lei. O efeito plausível é uma faixa: a
magnitude individual do análogo equatoriano, em três cenários que levam à
regra as ressalvas do Capítulo~\ref{cap:referencial}, multiplicada pela fração
de nascimentos exposta, em grade fixada no apêndice; decide a célula de 80
gramas e fração de 0,25.
Poder suficiente é condição necessária para o teste, não garantia de que o
efeito exista nessa magnitude. A data-limite para os registros é 31 de maio de
2027, fim do terceiro bimestre.

\begin{quadro}[htb]
\caption{\label{qua:regra}Regra de decisão}
\centering
\footnotesize
\begin{tabular}{@{}p{0.3cm}p{4.4cm}p{2.6cm}p{7.3cm}@{}}
\toprule
 & Medição dos voos & Poder & O que o ensaio entrega \\
\midrule
1 & Voos documentados por município e mês, com pico sazonal & Detectável $\leq$ plausível &
Desenho principal, confirmatório, com municípios e estação medidos. \\
2 & Idem & Detectável $>$ plausível &
Medida dos voos; desenho principal reportado como limite, sem teste. \\
3 & Voos documentados, sem pico sazonal & Qualquer &
Medida dos voos; o calendário não se aplica; forma reduzida descritiva. \\
4 & Série completa sem voos no período de referência & Qualquer &
Medida do alcance da lei; forma reduzida descritiva; sem estimativa de efeito da pulverização. \\
5 & Voos na referência e nenhum em 2018 (série completa): a prática cessou antes da lei & Qualquer &
Medida do alcance da lei; $\beta$ só descritivo, pois não distingue a lei do fim anterior da prática; sem estimativa de efeito da pulverização. \\
6 & Registros não obtidos, ou com menos de dois anos completos, até a data-limite & Com a classificação provisória &
Desenho principal com os municípios com aeronave no Censo de 2006 e a quadra chuvosa, exploratório (o parâmetro deste capítulo não é o alvo); cálculo de poder; medida parcial dos voos, se houver. \\
\bottomrule
\end{tabular}
\fonte{Elaboração própria. Definições e critérios no Apêndice~\ref{ap:regra}. A linha 4 prevalece sobre a 5, e a 5 sobre as linhas 1 a 3.}
\end{quadro}

\section{Inferência, poder e plano de análise}
\label{sec:inferencia}

Há poucos municípios tratados, contíguos, e a lei é atribuída no nível do
estado. Erros-padrão agrupados por município, com aproximação assintótica,
superestimam a precisão, e o \emph{wild cluster bootstrap} isolado não é
defensável com tão poucos tratados \cite{mackinnon2018wild}. O procedimento
principal é o de \citeonline{ferman2019inference}, que admite
heterocedasticidade por tamanho de município, aplicado à mudança do hiato
sazonal de cada município, líquida da chuva e dos efeitos fixos
(Apêndice~\ref{ap:regra}). Ele supõe municípios independentes, e os tratados
contíguos partilham choques; por isso o procedimento para um único agregado
tratado \cite{hagemann2025single} é reportado ao lado, e a conclusão
confirmatória exige que os dois concordem. Os de \citeonline{conley2011} e
\citeonline{canay2021wild} são sensibilidade, e a inferência conformal de
\citeonline{chernozhukov2021conformal} é usada no controle sintético.

Dois fatores limitam o poder e nenhum se resolve com mais anos de dados: o
número de municípios com aplicação aérea e a fração dos nascimentos de cada um
efetivamente exposta. O efeito mínimo detectável de $\beta$ tem nível de 5\%,
poder de 80\% e teste unilateral, sob o procedimento principal, e é calculado
apenas com dados do período de referência e do período de seca
(Apêndice~\ref{ap:regra}). Este projeto não traz o seu valor, o que é uma
limitação. A rodada exploratória de novembro, com classificação e estação
provisórias, indicou efeito detectável acima do plausível na célula que
decide, isto é, a linha 2 do Quadro~\ref{qua:regra}; a medição dos voos pode
mudar a classificação e o resultado. O valor pré-especificado é reportado no
plano de análise ou em adendo datado, antes da análise confirmatória.

\textbf{Rodadas exploratórias e plano de análise.} Duas rodadas exploratórias
antecedem este documento. Nenhuma estimativa delas é reportada; o
Apêndice~\ref{ap:regra} lista as variantes estimadas e o que estava fixado
antes de cada uma. A de setembro de 2026 estimou a forma reduzida estadual,
sob plano fechado no dia 21. A de novembro estimou o desenho de calendário,
ainda sem plano, com classificação e estação provisórias e sem a chuva, sobre
os nascimentos de 2005 a 2024: a exposição sazonal já foi cruzada com o
desfecho do pós-lei. Daí três compromissos. Municípios tratados e estação são
fixados pela medição dos voos, e os parâmetros ainda abertos, por critério
escrito sem uso do desfecho; a ratificação com o orientador não altera o
período de referência nem o Quadro~\ref{qua:regra}, e os limiares do
Apêndice~\ref{ap:regra} são provisórios até ela. Se municípios e estação medidos
coincidirem com variante já estimada, o resultado é rotulado como reestimação
em dados já examinados, e não como teste confirmatório. E o plano de análise
e o material das duas rodadas serão depositados em repositório público
versionado, com data verificável por terceiros, antes da análise
confirmatória, que é refeita a partir das bases originais.

% Seção 6 — Dados. Estado das fontes: CLAUDE.md (tabela de status) e
% docs/fontes-e-vintages.md. Sem descritivas de desfecho.

\chapter{Dados}
\label{cap:dados}

A unidade de análise é o município por mês de nascimento. O painel principal
cobrirá os municípios cearenses com aplicação aérea e os de comparação, de 2005
a 19 de dezembro de 2024; os 184 municípios do estado entram na forma reduzida,
e municípios de estados vizinhos, no controle sintético e na fronteira
(fontes no Apêndice~\ref{ap:fontes}).

\textbf{Desfechos.} Do SINASC serão obtidos o peso médio ao nascer, desfecho
confirmatório, e a proporção de nascidos com menos de 2.500 gramas e com menos
de 37 semanas de gestação, exploratórios (siglas no
Apêndice~\ref{ap:fontes}). A média é o desfecho confirmatório porque é nela
que o análogo mede o efeito e porque proporções de eventos raros são menos
precisas em células municipais pequenas; o baixo peso indica, sem teste, se o
efeito se concentra na cauda inferior. A série atravessa a troca do formulário de nascido vivo, em 2011
\cite{brasil2011sinasc}, que alterou o registro da idade gestacional; o peso ao
nascer não foi afetado do mesmo modo, e a prematuridade é analisada com essa
ressalva. Usa-se o município de residência da mãe. Do SIM será extraído o óbito
fetal, exploratório, e construído um desfecho composto, a proporção de
gestações registradas que terminam em nascido vivo com 2.500 gramas ou mais.
Do SINAN serão extraídas as notificações de intoxicação por agrotóxico
agrícola, por circunstância, para o canal de substituição
(Apêndice~\ref{ap:espec}), e as de arboviroses, para verificar se essa fonte
de sazonalidade do peso difere entre os grupos.

\textbf{Exposição.} A classificação pelo método vem dos registros da
Seção~\ref{sec:medicao} e, provisoriamente, do Censo Agropecuário de 2006; o
grupo de comparação, do Censo Agropecuário de 2017; a chuva, das séries de
precipitação por município. Como a lei estadual de 1993 autoriza os municípios
a legislar supletivamente, regras locais anteriores a 2019 alteram o período
de comparação; um levantamento preliminar encontrou uma proibição em Limoeiro
do Norte, de 20 de novembro de 2009 \cite{limoeiro2009lei1478}, revogada em 26
de maio de 2010 por lei que passou a regular a prática
\cite{limoeiro2010lei1511}. O levantamento será refeito, com conferência de
vigência.

\textbf{Escopo.} O canal hídrico, que exigiria medições de agrotóxicos na
água de consumo com localização dos pontos de coleta, não integra a análise
confirmatória, e a exposição não é medida no nível da residência. O §~2º do
art.~28-B alcança a dispersão aérea contra vetores; não se verificou se ela
ocorreu no Ceará antes de 2019, e, se ocorreu, há municípios tratados que a
classificação não vê.

\textbf{Reprodutibilidade e sigilo.} O processamento será feito por rotinas
versionadas, com data de extração registrada; microdados não são
redistribuídos, e células com poucos eventos não são divulgadas.

\textbf{Aspectos éticos.} Os dados de saúde são secundários, de acesso público
e sem identificação individual; pesquisas dessa natureza não são submetidas ao
sistema CEP/CONEP, nos termos da Resolução nº~510/2016 do Conselho Nacional de
Saúde \cite{cns2016res510}. Os registros obtidos pela Lei de Acesso à
Informação identificam empresas e responsáveis técnicos; serão usados apenas
para classificar municípios e meses, e nenhum resultado será divulgado com
identificação de empresa ou pessoa. Como os operadores são poucos, a medida
municipal e mensal pode identificá-los: agregação, limiar de supressão e
guarda dos documentos são fixados no plano de análise, antes do recebimento.

% Seção 7 — Ameaças à identificação. Flags 0–7 do CLAUDE.md + docs/ars/19 §6.

\chapter{Ameaças à identificação}
\label{cap:ameacas}

O Quadro~\ref{qua:ameacas} reúne as ameaças conhecidas, o que cada uma faz ao
parâmetro estimado e o que o desenho lhe opõe. Três merecem comentário.

\begin{quadro}[htb]
\caption{\label{qua:ameacas}Ameaças à identificação e respostas do desenho}
\centering
\footnotesize
\begin{tabular}{@{}p{3.5cm}p{4.9cm}p{6.1cm}@{}}
\toprule
Ameaça & Efeito sobre a estimativa & Resposta \\
\midrule
Tratamento medido pela cultura, não pelo método &
Atenuação proporcional ao $VPP$ (Eq.~\ref{eq:atenuacao}) &
Estágio 2 (medição) e estágio 3 (tratados pelo método). \\
Prática encerrada antes da lei (2011) &
Grupo tratado inexistente; nulo por ausência de tratamento &
Autos judiciais e relatórios do MAPA; regra de decisão (Quadro~\ref{qua:regra}). \\
Antecipação (2011, 2015 e 2018) &
Dose e desfecho respondem antes da vigência &
Sensibilidade com dose de 2010--2014; coortes de transição excluídas; \emph{leads} até 2018; o estágio 3 não usa a dose. \\
Grupo de comparação pequeno e contaminado &
Nível da estimativa depende de quem é o zero; controle vetorial aéreo também foi proibido &
Quatro construções do grupo; inferência adequada a poucos grupos; estágio 3 não depende desse grupo. \\
Seleção para o nascimento vivo &
Efeito sobre o peso médio atenuado ou invertido &
Óbito fetal como desfecho confirmatório. \\
Composição sazonal das mães &
Hiato sazonal reflete quem concebe em cada época &
Efeitos fixos de município por mês do ano; teste de estabilidade da composição materna. \\
Choques de 2020--2021 &
Mudança de sazonalidade de concepção e de assistência pré-natal &
Estimativas com e sem o período; doadores externos. \\
Substituição pelo aplicador terrestre &
Exposição populacional cai, ocupacional pode subir &
Notificações de intoxicação como canal; faz parte do efeito da política. \\
Fiscalização desconhecida &
Parâmetro é de intenção de tratar &
Declarado; registros de autuação solicitados nos pedidos de acesso. \\
\bottomrule
\end{tabular}
\fonte{Elaboração própria.}
\end{quadro}

\textbf{A medida do tratamento.} É a ameaça que organiza o projeto. A
proibição alcança o avião, e a variável disponível mede a lavoura. Se a
aplicação aérea era concentrada em poucos municípios, como indica o Censo
Agropecuário de 2006, a maior parte dos
classificados como tratados pela cultura não o era pelo método, e a
Equação~\ref{eq:atenuacao} implica que a forma reduzida pode ser pequena mesmo
que o efeito verdadeiro seja grande. Ampliar o grupo tratado pela cultura
tenderia a agravar o problema, ao acrescentar municípios sem registro de
avião. A resposta é medir o método e concentrar a
comparação onde ele existia.

\textbf{A persistência da prática.} Todas as estratégias que medem o método
dependem de a pulverização aérea ter continuado na Chapada do Apodi entre 2011
e 2018. O fato de a lei de 2019 ter sido aprovada contra uma prática descrita
como corrente sugere que sim, mas não o demonstra. Se os registros mostrarem
que ela cessou, o resultado do ensaio muda de natureza: a ausência de efeito
passa a ser ausência de tratamento, e a contribuição se desloca para a
mensuração do alcance efetivo da lei.

\textbf{A sazonalidade.} O desenho de calendário troca a hipótese de
tendências paralelas entre municípios pela de estabilidade relativa do hiato
sazonal. A troca é vantajosa se a aplicação aérea tiver de fato um pico
sazonal; se a aplicação for uniforme ao longo do ano, o desenho mede zero por
construção. Por isso a existência do pico é verificada nos relatórios mensais
antes da estimação, e o poder do desenho é calculado no pré-período.

% Capítulo 8 — Cronograma. As datas são PROPOSTA (seis bimestres a partir de
% dez/2026); ajustar ao calendário do programa.

\chapter{Cronograma}
\label{cap:cronograma}

O Quadro~\ref{qua:cronograma} distribui as atividades em
bimestres e cobre apenas o primeiro ensaio; a defesa segue o calendário do
programa. O plano de análise é depositado antes da análise confirmatória, e o
que for medido depois entra por adendo datado.

\begin{quadro}[htb]
\caption{\label{qua:cronograma}Cronograma de atividades, por bimestre}
\centering
\footnotesize
\begin{tabular}{@{}p{8.6cm}*{6}{>{\centering\arraybackslash}p{0.75cm}}@{}}
\toprule
Atividade & 1 & 2 & 3 & 4 & 5 & 6 \\
\midrule
Pedidos de acesso à informação e consulta à ação de 2011 & X & & & & & \\
Montagem do painel (2005--2024) e congelamento das bases & X & X & & & & \\
Plano de análise: ratificação com o orientador e depósito datado & & X & & & & \\
Diagnóstico do pré-período, poder e classificação dos municípios & & X & X & & & \\
Desenho principal, análises complementares e triangulação & & & X & X & & \\
Redação do ensaio e qualificação & & & & X & X & \\
Revisão final do ensaio & & & & & & X \\
\bottomrule
\end{tabular}
\fonte{Elaboração própria. Bimestre 1: dezembro de 2026 e janeiro de 2027.}
\end{quadro}


\postextual
\clearpage
\bibliography{../referencias,referencias-projeto}

% Apêndices (pós-textuais): A fontes, B ameaças, C especificações.
% Criados na revisão do Estágio 4 (2026-11-15), por decisão do pesquisador,
% para manter o corpo em 15 páginas.

\begin{apendicesenv}
\clearpage
\chapter{Fontes de dados}
\label{ap:fontes}

\begin{quadro}[htb]
\caption{\label{qua:fontes}Fontes de dados, uso e acesso}
\centering
\footnotesize
\begin{tabular}{@{}p{3.6cm}p{5.6cm}p{5.3cm}@{}}
\toprule
Fonte & Uso & Acesso, guarda e risco \\
\midrule
SINASC (DATASUS) & Peso ao nascer, baixo peso, prematuridade; 2005--2024 & Público. Troca de formulário em 2011 afeta a idade gestacional. \\
SIM (DATASUS) & Óbito fetal & Público. Subnotificação. \\
SINAN (DATASUS) & Intoxicações por agrotóxico; arboviroses & Público. \\
Censo Agropecuário 2006 (IBGE) & Estabelecimentos com aplicação por aeronave: classificação provisória & Público. Não distingue cultura. \\
Censo Agropecuário 2017 (IBGE) & Participação da agricultura familiar na banana: grupo de comparação & Público. \\
PAM (IBGE) & Área plantada por cultura: dose da forma reduzida & Público. \\
FAO-GAEZ & Aptidão agroclimática: grupo de dose nula da forma reduzida & Público. \\
Precipitação (FUNCEME, INMET) & Chuva da quadra chuvosa, por município e ano & Público. \\
SINASC de estados vizinhos & Doadores do controle sintético; fronteira & Público. \\
Relatórios mensais da aviação agrícola (MAPA) & Aplicação aérea por município e mês, 2008--2024, CE e RN & Pedido de acesso à informação. Em papel até 2021; preservação incerta. \\
Cadastro de prestadores e autos de infração (órgãos estaduais) & Operadores; cumprimento após a lei & Pedido de acesso à informação. \\
Receituário agronômico & Modalidade de aplicação prescrita & Guarda a localizar; registra prescrição. \\
Autos judiciais (Justiça Federal) & Prática na Chapada após 2011 & Consulta processual. \\
Leis municipais & Proibições e regras locais anteriores a 2019 & Público (câmaras municipais). \\
\bottomrule
\end{tabular}
\fonte{Elaboração própria.}
\end{quadro}

\clearpage
\chapter{Ameaças à identificação e respostas do desenho}
\label{ap:ameacas}

\begin{quadro}[htb]
\caption{\label{qua:ameacas}Ameaças à identificação e respostas do desenho}
\centering
\footnotesize
\begin{tabular}{@{}p{3.6cm}p{4.7cm}p{6.2cm}@{}}
\toprule
Ameaça & Efeito sobre a estimativa & Resposta \\
\midrule
Voos suspensos pela seca antes da lei &
A lei impede a volta, não interrompe; efeito depende da chuva &
Referência em 2007--2011; período de seca em separado; relatórios mensais. \\
Seca com efeito próprio e fim próximo da lei &
Confunde a lei com o retorno das chuvas &
Chuva como controle nos dois grupos; especificação com efeito distinto por grupo; limitação declarada. \\
Retração da lavoura irrigada e do emprego &
Muda renda e composição materna nos tratados &
Efeitos fixos de município por ano; características maternas como desfecho-placebo. \\
Diluição dentro do município &
Efeito individual multiplicado pela fração exposta &
Poder calculado para uma grade de frações; resultado sem poder lido como limite. \\
Mudança de sazonalidade só nos tratados, após a lei &
Viola a hipótese de identificação &
Estudo de evento do hiato sazonal; tratamento falso; estimativas com e sem 2020--2021. \\
Classificação provisória de 2006 &
Tratados mal identificados; atenuação &
Relatórios mensais; exclusão de qualquer município com aeronave do grupo de comparação. \\
Seleção para o nascimento vivo &
Efeito sobre o peso médio atenuado ou invertido &
Óbito fetal e desfecho composto, exploratórios. \\
Interferência entre municípios &
Sinal indeterminado &
Três grupos; doadores fora do alcance; registros do estado vizinho. \\
Substituição terrestre, drones e descumprimento &
Parâmetro é de intenção de tratar &
Registros posteriores à lei; notificações de intoxicação. \\
Antecipação (projetos de lei de 2011 e 2015) &
Comportamento muda antes da vigência &
Referência anterior aos projetos; coortes de transição excluídas. \\
\bottomrule
\end{tabular}
\fonte{Elaboração própria.}
\end{quadro}

\clearpage
\chapter{Especificações}
\label{ap:espec}

\begin{quadro}[htb]
\caption{\label{qua:espec}Especificações do desenho de calendário}
\centering
\footnotesize
\begin{tabular}{@{}p{3.4cm}p{6.0cm}p{5.1cm}@{}}
\toprule
Análise & Especificação & Amostra e estatuto \\
\midrule
Principal ($H_1$) & Equação~\ref{eq:ddd}; coeficiente $\beta$ & Municípios com aplicação aérea e de banana familiar; confirmatória. \\
Período de seca ($H_2$) & Equação~\ref{eq:ddd}; coeficiente $\beta_S$ & Mesma amostra; descritiva. \\
Estudo de evento & Um coeficiente de $A_m \times P_t$ por ano, com 2007--2011 como base & Mesma amostra; verificação da hipótese de identificação. \\
Chuva como controle & Acrescenta $P_t \times C_{m,a}$, com $C_{m,a}$ a chuva da quadra no município e ano & Mesma amostra; parte da especificação principal. \\
Chuva por grupo & Acrescenta $A_m \times P_t \times C_{m,a}$ & Mesma amostra; sensibilidade. \\
Efeito por chuva ($H_3$) & Acrescenta $A_m \times P_t \times Pos_t \times C_{m,a}$ & Mesma amostra; exploratória. \\
Terceiro trimestre & $P_t$ redefinido como fração do terceiro trimestre na estação & Mesma amostra; exploratória. \\
Tratamento falso ($H_4$) & Equação~\ref{eq:ddd} com banana familiar no lugar de $A_m$ & Municípios de banana familiar e sem banana; falsificação. \\
Interferência & Equação~\ref{eq:ddd} com indicador de vizinho imediato de município tratado & Vizinhos e demais não tratados; exploratória. \\
Composição materna & Idade e escolaridade da mãe como desfecho na Equação~\ref{eq:ddd} & Mesma amostra; placebo. \\
\bottomrule
\end{tabular}
\fonte{Elaboração própria. A quadra chuvosa é tomada, provisoriamente, como
fevereiro a maio; os meses e a medida de chuva são fixados no plano de
análise.}
\end{quadro}
\end{apendicesenv}


\end{document}
